% Generated draft: requires mathematical and source review.
\documentclass[11pt]{article}
\usepackage[margin=25mm]{geometry}
\usepackage{fontspec}
\setmainfont{Latin Modern Roman}
\usepackage{amsmath,amssymb,mathtools,unicode-math}
\setmathfont{Latin Modern Math}
\usepackage{enumitem,needspace,xurl,hyperref}
\setlength{\parindent}{0pt}
\setlength{\parskip}{5pt}
\setlength{\emergencystretch}{3em}
\newcommand{\N}{\mathbb{N}}
\newcommand{\Z}{\mathbb{Z}}
\newcommand{\Q}{\mathbb{Q}}
\newcommand{\R}{\mathbb{R}}
\newcommand{\C}{\mathbb{C}}
\newcommand{\F}{\mathbb{F}}
\newcommand{\A}{\mathbb{A}}
\newcommand{\PP}{\mathbb{P}}
\newcommand{\E}{\mathbb{E}}
\newcommand{\eps}{\varepsilon}
\newcommand{\dd}{\,\mathrm{d}}
\begin{document}
\section*{UnsolvedMath ID: 20001910}
\textbf{Source problem code:} catalogue entry
\subsection{Definitions and mathematical statement}
The source asks whether a closed manifold admits a Riemannian metric with everywhere positive sectional curvature and ergodic geodesic flow, especially a metric close to round on \(S^2\). The aid uses connected closed manifolds of dimension \(n\ge2\) and normalized Liouville probability \(\mu_L\) on the unit tangent bundle. Ergodicity means every invariant measurable set has measure zero or one; it is not mere topological transitivity. The topology of near-round closeness is unspecified. The aid reports the existence question apparently open through 14 August 2026.

The reported necessary-condition package starts with \(\kappa_{min}=\min\sec_g>0\). It gives finite fundamental group and, along every geodesic, a conjugate point to its initial point by length \(\pi/\sqrt{\kappa_{min}}\), hence excludes an Anosov Riemannian geodesic flow under the cited no-conjugate-points input. Positive curvature alone is not said to imply ergodicity.

For an exact quantitative symmetry obstruction, let \(X\) be a nonzero Killing field and set
\[
 A=\max_M|X|,\quad U=\{p:|X(p)|>A/2\},\quad
 c_n=\frac{\sigma_{n-1}(\{u\in S^{n-1}:u_1>1/2\})}
             {\sigma_{n-1}(S^{n-1})}>0.
\]
Its conserved geodesic momentum \(J_X(p,v)=g(X(p),v)\) gives disjoint invariant sets
\[
 E_+=\{J_X>A/4\},\qquad E_-=\{J_X<-A/4\},
\]
and the aid reports
\[
 \mu_L(E_+),\ \mu_L(E_-)
 \ge c_n\frac{\operatorname{Vol}_g(U)}{\operatorname{Vol}_g(M)}>0.
\]
Thus any Liouville-ergodic candidate has no nonzero Killing vector field. More generally, it reports that a symmetric Killing tensor, whose homogeneous polynomial in velocities is a geodesic first integral, must vanish in odd degree and in even degree \(2k\) be a constant energy power, equivalently a constant multiple of \(g^{\odot k}\) under the chosen tensor convention. Consequently the isometry group is finite; this does not exclude a discrete finite symmetry group.

For a positively curved \(S^2\), the aid further reports positive topological entropy and a hyperbolic invariant set as necessary consequences of ergodicity. In the named \(9/16\)-pinched near-round regime it reports a required parabolic or degenerate elliptic closed geodesic. The exact entropy theorem, pinching strictness, and linear-versus-nonlinear degeneracy convention are absent; this last diagnostic is therefore retained as an incomplete named condition, not a fully specified test or a sufficient construction.

The appended Section 8.3 is explicitly identified by the source verification as neighboring context for Question 8.3.1: unit vectors are grouped by whether lifted geodesics stay bounded (\(T_b\)), escape properly (\(T_p\)), or neither (\(T_i\)); the text says these are measurable and invariant, so one has full measure under ergodicity. The extraction inconsistently writes \(TM\) for the complement ambient space and does not specify the escape direction. Those sets are not hypotheses of Question 8.2.1 and are not used in the obstruction above. No positive-curvature ergodic example or universal nonexistence theorem is supplied.
\subsection{Short English statement}
The aid reports necessary gates for positive-curvature Liouville ergodicity, including a quantitative Killing-field obstruction and a restriction on all Killing tensors. Additional surface/near-round claims lack their exact theorem conventions. The existence question is reported apparently open at the stated August 2026 cutoff.
\subsection{Sources}
\begin{itemize}
\item Local UnsolvedMath registry, permanent ID \(20001910\), source code \texttt{catalogue entry}; \texttt{aim-geometry-notes.json}, \texttt{notes[247]}, tag \texttt{question}.
\item AIM workshop ``Geodesics,'' recorded item 8.2, extracted Question 8.2.1: \url{https://aimath.org/pastworkshops/geodesicsproblems.pdf}. Source type PDF; statement status \texttt{exact}. Supplied verification explicitly identifies the following \(T_b,T_p,T_i\) definitions as Section 8.3 contamination, preserved in the source snapshot rather than used as hypotheses. No PDF was accessed or compiled.
\item No authoritative citations, theorem identifiers, exact surface/pinching hypotheses, or full proof artifacts accompany the focusing, entropy, degeneracy, or Killing-tensor claims. Original remarks and literature fields are empty.
\item Research attempt 2: \texttt{partial\_result}, status at run \texttt{open}, rigor label \texttt{proved}, novelty confidence low; registry status \texttt{partially\_solved}, classification \texttt{PARTIAL-PROGRESS}. Created and updated 14 August 2026; the aid says apparently open through 2026-08-14. Machine-generated aid, not peer reviewed; independently verify before citation.
\item Ulam AI / UnsolvedMath Contributors; CC BY 4.0.
\end{itemize}
\end{document}
